\enquestion{4.1}{
    Suppose that $f(x) = 2e^{-2x}$ for $0 < x$. Determine the following:
    \begin{center}
        \begin{tabular}{*{3}{m{0.3\textwidth}}}
            (a) $P(1 < X)$ & (b) $P(1 < X < 2.5)$ & \multirow{2}{*}{(e) $P(3 \leq X)$}\\
            (c) $P(X = 3)$ & (d) $P(X < 4)$ & \\
            \multicolumn{2}{l}{(f) Determine $x$ such that $P(x < X) = 0.20$.} \\
            \multicolumn{2}{l}{(g) Determine $x$ such that $P(X \le x) = 0.10$.}
        \end{tabular}
    \end{center}
}
\enanswer{
    This function is a probability density function of a continuous random variable $X$ (specifically, the exponential distribution).
    We can compute probabilities on intervals by using integrals and probabilities for specific outcomes are always $0$.
    \begin{enumerate}[label=(\alph*)]
        \item $P(1 < X) = \int_{x=1}^{\infty} 2e^{-2x} \dx \approx 0.1353$,
        \item $P(1 < X < 2.5) = \int_{x=1}^{2.5} 2e^{-2x} \dx \approx 0.1286$,
        \item $P(X = 3) = 0$ (as $X$ is a continuous random variable!),
        \item $P(X < 4) = \int_{x=0}^4 2e^{-2x} \dx \approx 0.9997$,
        \item $P(3 \leq X) = \int_{x=3}^{\infty} 2e^{-2x} \dx \approx 0.00025$,
        \item To determine $x$ such that $P(x < X) = 0.20$ (or $P(X \leq x) = 0.80$), we need to use the \enquote{numeric solver}:
        \begin{align*}
            E_1 &= \int_{0}^{x} 2e^{-2x} \dx, \\
            E_2 &= 0.8.
        \end{align*}
        The numeric solver gives us $x \approx 0.8047$,
        \item To determine $x$ such that $P(X \le x) = 0.10$, we need to use the \enquote{numeric solver}:
        \begin{align*}
            E_1 &= \int_{0}^{x} 2e^{-2x} \dx, \\
            E_2 &= 0.1.
        \end{align*}
        The numeric solver gives us $x \approx 0.0527$.
    \end{enumerate}
}

\enquestion{4.7}{
    Suppose that $f(x) = \frac{3}{2}x^2$ for $-1 < x < 1$. Determine the following:

    \begin{center}
        \begin{tabular}{*{2}{m{0.45\textwidth}}}
            (a) $P(0 < X)$ & (b) $P(0.5 < X)$ \\
            (c) $P(-0.5 \le X \le 0.5)$ & (d) $P(X < -2)$ \\
            (e) $P(X < 0\text{ or } X > -0.5)$ & (f) $x$ such that $P(x < X) = 0.05$.
        \end{tabular}
    \end{center}
}
\enanswer{
    This function is a probability density function of a continuous random variable $X$.
    We can compute probabilities on intervals by using integrals and probabilities for specific outcomes are always $0$.
    \begin{enumerate}[label=(\alph*)]
        \item $P(0 < X) = \int_{x=-1}^{0} \frac{3}{2}x^2 \dx = 0.5$ (notice that the distribution is symmetric around $x = 0$, so we did not need to compute this integral to see this),
        \item $P(0.5 < X) = \int_{x=0.5}^{1} \frac{3}{2}x^2 \dx = 0.4375$,
        \item $P(-0.5 \le X \le 0.5) = \int_{x=-0.5}^{0.5} \frac{3}{2}x^2 \dx = 0.125$,
        \item $P(X < -2) = 0$ (all of the density is between $-1$ and $1$, so this probability is $0$),
        \item $P(X < 0\text{ or } X > -0.5) = 1$ (notice that all possible outcomes are either smaller than zero, or larger than $-0.5$),
        \item To determine $x$ such that $P(x < X) = 0.05$ (or $P(X \leq x) = 0.95$), we need to use the \enquote{numeric solver}:
        \begin{align*}
            E_1 &= \int_{-1}^{x} \frac{3}{2}x^2 \dx, \\
            E_2 &= 0.95.
        \end{align*}
        The numeric solver gives us $x \approx 0.9655$.
    \end{enumerate}
}

\enquestion{4.11}{
    The probability density function of the length of a metal rod is $f(x) = 2$ for $2.3 < x < 2.8$ meters.
    \begin{enumerate}[label=(\alph*)]
        \item If the specifications for this process are from $2.25$ to $2.75$ meters, what proportion of rods fail to meet the specifications?
        \item Assume that the probability density function is $f(x) = 2$ for an interval of length $0.5$ meters. 
        Over what value should the density be centered to achieve the greatest proportion of rods within specifications?
    \end{enumerate}
}
\enanswer{
    \begin{enumerate}[label=(\alph*)]
        \item Notice that this is the probability density function of the continuous uniform distribution between $a=2.3$ and $b=2.8$.
        The rods that fail to meet the specifications are either (i) smaller than $2.25$ meter or (ii) longer than $2.75$ meters.
        Only the second case is possible.
        The proportion of rods failing to meet the specifications is
        \[
            P(X \geq 2.75) = \int_{x=2.75}^{2.8} 2 \dx = 2 \cdot 0.05 = 0.1. 
        \] 

        \item All rods would achieve the specifications if the distribution would be from $2.25$ to $2.75$ (this is indeed an interval of length $0.5$ meters).
        In other words, the density should be centered at $\frac{2.25 + 2.75}{2} = 2.5$.
    \end{enumerate}
}

\enquestion{4.13}{
    Suppose that the cumulative distribution function of the random variable $X$ is
    \[
        F(x) = \begin{cases}
            0 & x < 0 \\
            0.2x & 0 \le x < 5 \\
            1 & 5 \le x
        \end{cases}
    \]
    Determine the following:

    \begin{center}
        \begin{tabular}{*{2}{m{0.45\textwidth}}}
            (a) $P(X < 2.8)$ & (b) $P(X > 1.5)$ \\
            (c) $P(X < -2)$ & (d) $P(X > 6)$
        \end{tabular}
    \end{center}
}
\enanswer{
    We can use the specified CDF to read the desired probabilities:
    \begin{enumerate}[label=(\alph*)]
        \item $P(X < 2.8) = P(X \leq 2.8) = F(2.8) = 0.2 \cdot 2.8 = 0.56$,
        \item $P(X > 1.5) = 1 - P(X \leq 1.5) = 1 - F(1.5) = 1 - 0.2 \cdot 1.5 = 0.7$,
        \item $P(X < -2) = F(-2) = 0$,
        \item $P(X > 6) = 1 - P(X \leq 6) = 1 - F(6) = 1 - 1 = 0$.
    \end{enumerate}
}

\enquestion{4.15}{
    Determine the cumulative distribution function given the following probability density function:
    \[
        f(x) = 2e^{-2x} \qquad \text{ for } 0 < x
    \]
}
\enanswer{
    Notice that $f(x) = 0$ for all $x \leq 0$, so is $F(x) = 0$ for all $x \leq 0$.
    For $x > 0$, notice that the cumulative distribution function is the antiderivative of the probability density function:
    \begin{align*}
        F(x) = P(X \leq x) = \int_{y=0}^{x} f(y) \dy &= \int_{y=0}^{x} 2e^{-2y} \dy \\
                                                     &= \left[ -e^{-2y} \right|_{y=0}^x \\
                                                     &= -e^{-2x} - (-e^{-2\cdot 0}) \\
                                                     &= 1 - e^{-2x}.
    \end{align*}
    Therefore, the cumulative distribution function is equal to
    \[
        F(x) = \begin{cases}
            0 & x \le 0 \\
            1 - e^{-2x} & x \ge 0
        \end{cases}
    \]
}

\enquestion{4.21}{
    Determine the cumulative distribution function for the following probability density function:
    \[
        f(x) = \begin{cases}
                    2, & \text{ for } 2.3 < x < 2.8 \\
                    0, & \text{ otherwise.}
        \end{cases}
    \]
    Use the cumulative distribution function to determine the probability that a length exceeds $2.7$ meters.
}
\enanswer{
    Notice that $f(x) = 0$ for all $x \leq 2.3$, so $F(x) = 0$ for all $x \leq 2.3$.
    For $2.3 < x < 2.8$, the cumulative distribution function is the antiderivative of the probability density function:
    \begin{align*}
        F(x) = P(X \leq x) = \int_{y=2.3}^{x} f(y) \dy &= \int_{y=2.3}^{x} 2 \dy \\
                                                     &= \left[ 2y \right|_{y=2.3}^x \\
                                                     &= 2 \cdot (x - 2.3).
    \end{align*}
    For $x \geq 2.8$, $f(x) = 0$, so we have reached the peak of $F(x) = 1$ for $x \geq 2.8$.

    Therefore, the cumulative distribution function is equal to
    \[
        F(x) = \begin{cases}
            0 & x \le 2.3 \\
            2\cdot (x - 2.3) & 2.3 \le x \le 2.8 \\
            1 & x \ge 2.8
        \end{cases}
    \]
    The probability that a length exceeds $2.7$ meters is then equal to 
    \[
        P(X > 2.7) = 1 - P(X \leq 2.7) = 1 - F(2.7) = 1 - 2 \cdot (2.7 - 2.3) = 0.2.
    \]
}

\enquestion{4.29}{
    Suppose that $f(x) = \frac{3}{2}x^2$ for $-1 < x < 1$. 
    Determine the mean and variance of $X$.
}
\enanswer{
    Using the general formulas for expected value and variance of a random variable $X$, we get:
    \begin{align*}
        E(X) = \int_{x=-\infty}^{\infty} x \cdot f(x) \dx &= \int_{x=-1}^{1} x \cdot \frac{3}{2}x^2 \dx \\
                                                        &= \int_{x=-1}^{1} \cdot \frac{3}{2}x^3 \dx \\
                                                        &= \left[ \frac{3}{8}x^4\right|_{x=-1}^1 \\
                                                        &= \frac{3}{8} - \frac{3}{8} \\
                                                        &= 0.
    \end{align*}
    \begin{align*}
        V(X) = \int_{x=-\infty}^{\infty} (x - E(X))^2 \cdot f(x) \dx &= \int_{x=-1}^{1} (x - 0)^2  \cdot \frac{3}{2}x^2 \dx \\
                                                        &= \int_{x=-1}^{1} \cdot \frac{3}{2}x^4 \dx \\
                                                        &= \left[ \frac{3}{10}x^5\right|_{x=-1}^1 \\
                                                        &= \frac{3}{10} - (-\frac{3}{10}) \\
                                                        &= \frac{6}{10} = 0.6.
    \end{align*}
}

\enquestion{4.37}{
    Integration by parts is required.
    The probability density function of the diameter of a drilled hole in millimeters is $f(x) = 10 e^{-10(x-5)}$ for $x > 5$ mm.
    Although the target diameter is $5$ millimeters, vibrations, tool wear, and other nuisances produce diameters greater than $5$ millimeters.
    \begin{enumerate}[label=(\alph*)]
        \item Determine the mean and variance of the diameter of the holes.
        \item Determine the probability that a diameter exceeds $5.1$ millimeters.
    \end{enumerate}
}
\enanswer{
    Using the general formulas for expected value and variance of a random variable $X$, we get:
    \begin{align*}
        E(X) &=\int_{x=-\infty}^{\infty} x \cdot f(x) \dx \\
             &=\int_{x=5}^{\infty} 10 x e^{-10(x-5)} \dx.
    \end{align*}

    Using integration by parts, with
    \[
        u=x, \qquad dv=10e^{-10(x-5)}\dx,
    \]
    we obtain
    \[
        du=dx, \qquad v=-e^{-10(x-5)}.
    \]

    Hence,
    \begin{align*}
        E(X)    &= \left[-xe^{-10(x-5)} \right]_{x=5}^{\infty} - \int_{x=5}^{\infty} -e^{-10(x-5)} \dx \\
                &= 5+\left[-\frac{1}{10}e^{-10(x-5)}\right]_{x=5}^{\infty}\\
                &= 5+\frac{1}{10} = 5.1.
    \end{align*}


    We use the identity
    \[
        V(X) = E(X^2) - (E(X))^2.
    \]

    First, calculate the second moment:
    \[
        E(X^2)=\int_5^\infty 10x^2e^{-10(x-5)}\dx.
    \]

    Applying integration by parts with
    \[
        u=x^2, \qquad dv=10e^{-10(x-5)}\dx,
    \]
    gives
    \[
        du=2x\dx, \qquad v=-e^{-10(x-5)}.
    \]

    Consequently,
    \begin{align*}
    E(X^2)  &=\left[-x^2e^{-10(x-5)}\right]_{x=5}^{\infty} +2\int_{x=5}^{\infty} xe^{-10(x-5)}\dx\\
            &=25+\frac{2}{10}E(X)\\
            &=25+\frac{2}{10}(5.1)\\
            &=26.02.
    \end{align*}

    Thus,
    \begin{align*}
    V(X)    &=E(X^2)-(E(X))^2\\
            &=26.02-(5.1)^2\\
            &=26.02-26.01 = 0.01.
    \end{align*}

    We calculate the probability by integrating the density:
    \begin{align*}
    P(X>5.1)    &=\int_{x=5.1}^\infty 10e^{-10(x-5)}\dx\\
                &=\left[-e^{-10(x-5)}\right|_{x=5.1}^{\infty}\\
                &=e^{-10(5.1-5)}\\
                &=e^{-1} \approx 0.3679.
    \end{align*}
}

\enquestion{4.41}{
    The thickness of a flange on an aircraft component is uniformly distributed between $0.90$ and $1.05$ millimeters.
    \begin{enumerate}[label=(\alph*)]
        \item Determine the cumulative distribution function of flange thickness.
        \item Determine the proportion of flanges that exceeds 1.02 millimeters.
        \item What thickness is exceeded by \SI{90}{\percent} of the flanges?
        \item Determine the mean and variance of flange thickness.
    \end{enumerate}
}
\enanswer{
    The cumulative distribution function of a continuous uniform distribution from $a$ to $b$ is equal to
    \[
        F(x) = \begin{cases} 0, & \text{ for } x < a \\ \frac{x - a}{b - a}, & \text{ for } a \leq x < b \\ 1, & \text{ for } x \geq b \end{cases}
    \]
    \begin{enumerate}[label=(\alph*)]
        \item With other words, the cumulative distribution function for $a = 0.90$ and $b = 1.05$ equals
        \[
            F(x) = \begin{cases} 0, & \text{ for } x < 0.90 \\ \frac{x - 0.90}{0.15}, & \text{ for } 0.90 \leq x < 1.05 \\ 1, & \text{ for } x \geq 1.05 \end{cases}
        \]
        \item $P(X > 1.02) = 1 - P(X \leq 1.02) = 1 - F(1.02) = 1 - \frac{1.02 - 0.90}{0.15} = 0.2$,
        \item The thickness exceeded by \SI{90}{\percent} of the flanges is at a tenth of the interval from $0.90$ to $1.05$.
        With other words, this thickness is equal to $a + 0.1(b - a) = 0.90 + 0.1(1.05 - 0.9) = 0.915$. 
        \item The mean and variance of a uniform distribution are equal to $\frac{a+b}{2}$ and $\frac{(b-a)^2}{12}$ respectively.
        Therefore, we have:
        \begin{align*}
            E(X) &= \frac{0.9 + 1.05}{2} = 0.975,\\
            V(X) &= \frac{(1.05 - 0.9)^2}{12} = 0.001875.
        \end{align*}
    \end{enumerate}
}